\section*{Reparto del trabajo}
\addcontentsline{toc}{section}{Reparto del trabajo}

El trabajo se repartió por ficheros con contratos fijos entre ambos (\code{src/common.py}), de
modo que cada autor pudiera avanzar con datos o modelos simulados sin esperar al otro.

\begin{itemize}
  \item \textbf{Daniel Rodríguez Alonso}: descarga y verificación de las anotaciones
    (\code{download.py}, \code{show\_boxes.py}), recorte de manos (\code{prepare\_data.py}),
    equilibrado y partición por persona (\code{make\_splits.py}), carga de datos y aumento
    (\code{dataset.py}), evaluación y curvas (\code{evaluate.py}, \code{plot\_curves.py}),
    Grad-CAM (\code{gradcam.py}), análisis estadístico y de coste (\code{analysis.py}) y ejecución
    de los entrenamientos de la ablación y de la referencia.
  \item \textbf{Carlos Falcón Brito}: propuesta inicial del proyecto, arquitectura GestureCNN y
    ResNet18 (\code{model.py}), datos simulados para desarrollar sin el conjunto real
    (\code{mock\_data.py}), bucle de entrenamiento con parada anticipada (\code{train.py}),
    configuraciones de la ablación y lanzador de experimentos (\code{configs/},
    \code{run\_experiments.py}) y demo con ROS\,2 (\cref{sec:demo}).
\end{itemize}

\paragraph{Uso de herramientas de IA.}
Se utilizó un asistente de IA (Claude, de Anthropic) como apoyo en la programación, en la
ejecución de los experimentos y en la redacción de este informe. Todas las cifras proceden de los
ficheros de resultados del repositorio, generados por los \emph{scripts} descritos. La
responsabilidad sobre el contenido corresponde a los autores.

\paragraph{Código.}
El código, las configuraciones, la partición y los resultados están en
\url{https://github.com/danirodriguezz/gestos-cnn-robot}. Las cifras de este documento se generan
automáticamente a partir de \code{results/} con \code{src/analysis.py}.
